El Premi Nobel de Física de l'any 2020 es va concedir a tres científics: Andrea Ghez, Reinhard Genzel i Roger Penrose (Ghez és la quarta dona de la història que aconsegueix un Nobel de Física). En concret, Ghez i Genzel van ser guardonats per les seves deduccions sobre Sagitari A* (el forat negre supermassiu del centre de la Via Làctia) a partir de l'observació de les òrbites d'estels propers. La gràfica adjunta mostra la intensitat del camp gravitatori provocat per Sagitari A* en funció de la distància.

\figura[0.55]{fig1.png}

\begin{apartat}{1,25}
Al voltant de Sagitari A* hi orbita una estrella anomenada S2. Quin és el mòdul de la velocitat d'aquesta estrella quan es troba a una distància de $3,00\times 10^{13}\ \mathrm{m}$?
\end{apartat}

\begin{apartat}{1,25}
Calculeu la massa de Sagitari A*. A quantes masses del Sol equival? Com que Sagitari A* és un forat negre, fins a una certa distància d'ell la velocitat d'escapament de Sagitari A* és superior a la velocitat de la llum. Calculeu a partir de quina distància aquesta velocitat d'escapament és menor que la velocitat de la llum. Cal deduir la velocitat a partir de consideracions energètiques.
\end{apartat}

\nota{Podeu negligir els efectes relativistes i podeu suposar que la llei de la gravitació universal és vàlida.}
\dades{%
  $G = 6,67\times 10^{-11}\ \mathrm{N\,m^2\,kg^{-2}}$.\\
  $M_\mathrm{Sol} = 1,99\times 10^{30}\ \mathrm{kg}$.\\
  $c = 3,00\times 10^{8}\ \mathrm{m\,s^{-1}}$.}
